% GENERATED FILE. Edit canonical lessons, then run npm run generate:books.
% canonical-source-hash: fnv1a64:66a8d29b09a5004f
% canonical-lessons: GU-C494-balpan, GU-C494-ucharvu, GU-C494-kisor, GU-C494-vrddh, GU-C494-vrddhasram

\chapter{Growing Up, Growing Old}
\label{ch:gu-growing-up-growing-old}
\hlchaptermodality{494}

\begin{chapteropening}
\textbf{By the end of this chapter:} \emph{I can talk about childhood and old age.}

Complete the last lesson of chapter 494: I can talk about childhood and old age.
\end{chapteropening}

\section[\texorpdfstring{\gu{બાળપણ} (bāḷpaṇ)}{બાળપણ (bāḷpaṇ)}]{\gu{બાળપણ} (bāḷpaṇ) — childhood}

\label{lesson:GU-C494-balpan}



\begin{quote}
Before the new one: say the Gujarati for a household, then the Gujarati for offspring.
\end{quote}

\subsection*{You'll want to know: \gu{બાળપણ}}
\textbf{\gu{બાળપણ}} — \emph{bāḷpaṇ} — \textquotedblleft{}childhood\textquotedblright{}.

\begin{quote}
\gu{દાદા} \gu{પોતાના} \gu{બાળપણની} \gu{વાર્તા} \gu{કહે} \gu{છે}.
\end{quote}

\begin{quote}
\emph{dādā potānā bāḷpaṇnī vārtā kahe che.}
\end{quote}

\begin{quote}
Grandfather tells stories of his childhood.
\end{quote}

A neuter noun.

\begin{cousinweb}
\textbf{\gu{બાળપણ}} (\emph{bāḷpaṇ}) goes back to a reconstructed \emph{*bālatvana-}, \textquotedblleft{}childhood\textquotedblright{}, built on Sanskrit \emph{bāla}, \textquotedblleft{}a child\textquotedblright{}; CDIAL lists the Gujarati word (CDIAL 9218).
\end{cousinweb}

\begin{grammarlens}[title={What you've built}]
One more word for describing what you lived and felt.
\end{grammarlens}

\subsection*{Your turn}
\begin{itemize}
\raggedright
  \item \emph{Say it:} \emph{bāḷpaṇ}
  \item \emph{Say it:} \emph{bāḷpaṇ}, once more
  \item \emph{Say it:} \emph{dādā potānā bāḷpaṇnī vārtā kahe che.} — Grandfather tells stories of his childhood.
  \item \emph{Recall:} say the Gujarati for a household, then the Gujarati for offspring, then say \emph{bāḷpaṇ} again
\end{itemize}

\begin{tcolorbox}[breakable,colback=teal!4,colframe=teal!35!black,title={Before you move on}]
What does \gu{બાળપણ} mean? (\textquotedblleft{}Childhood\textquotedblright{}.) Say it once more.
\end{tcolorbox}
\section[\texorpdfstring{\gu{ઉછરવું} (ucharvũ)}{ઉછરવું (ucharvũ)}]{\gu{ઉછરવું} (ucharvũ) — to grow up; to be brought up}

\label{lesson:GU-C494-ucharvu}



\begin{quote}
Before the new one: say the Gujarati for offspring, then the Gujarati for childhood.
\end{quote}

\subsection*{You'll want to know: \gu{ઉછરવું}}
\textbf{\gu{ઉછરવું}} — \emph{ucharvũ} — \textquotedblleft{}to grow up; to be brought up\textquotedblright{}.

\begin{quote}
\gu{મીરા} \gu{ગામમાં} \gu{ઉછરી} \gu{છે}.
\end{quote}

\begin{quote}
\emph{Mīrā gāmmā̃ uchrī che.}
\end{quote}

\begin{quote}
Mira grew up in a village.
\end{quote}

Past (he …): \gu{તે} \gu{મામાના} \gu{ઘરે} \gu{ઉછર્યો}, \emph{te māmānā ghare ucharyo}.

\begin{cousinweb}
\textbf{\gu{ઉછરવું}} (\emph{ucharvũ}) is inherited from Sanskrit \emph{ucchrīyatē}, \textquotedblleft{}is raised\textquotedblright{}; CDIAL lists the Gujarati verb, \textquotedblleft{}to be brought up\textquotedblright{} (CDIAL 1864). The causative \textbf{\gu{ઉછેરવું}} (\emph{uchervũ}), \textquotedblleft{}to bring up\textquotedblright{}, is its partner.
\end{cousinweb}

\begin{grammarlens}[title={What you've built}]
One more word for describing what you lived and felt.
\end{grammarlens}

\subsection*{Your turn}
\begin{itemize}
\raggedright
  \item \emph{Say it:} \emph{ucharvũ}
  \item \emph{Say it:} \emph{ucharvũ}, once more
  \item \emph{Say it:} \emph{Mīrā gāmmā̃ uchrī che.} — Mira grew up in a village.
  \item \emph{Recall:} say the Gujarati for offspring, then the Gujarati for childhood, then say \emph{ucharvũ} again
\end{itemize}

\begin{tcolorbox}[breakable,colback=teal!4,colframe=teal!35!black,title={Before you move on}]
What does \gu{ઉછરવું} mean? (\textquotedblleft{}To grow up; to be brought up\textquotedblright{}.) Say it once more.
\end{tcolorbox}
\section[\texorpdfstring{\gu{કિશોર} (kiśor)}{કિશોર (kiśor)}]{\gu{કિશોર} (kiśor) — a teenage boy}

\label{lesson:GU-C494-kisor}



\begin{quote}
Before the new one: say the Gujarati for childhood, then the Gujarati for to grow up.
\end{quote}

\subsection*{You'll want to know: \gu{કિશોર}}
\textbf{\gu{કિશોર}} — \emph{kiśor} — \textquotedblleft{}a teenage boy\textquotedblright{}.

\begin{quote}
\gu{તે} \gu{કિશોર} \gu{ઘણો} \gu{હોશિયાર} \gu{છે}.
\end{quote}

\begin{quote}
\emph{te kiśor ghaṇo hośiyār che.}
\end{quote}

\begin{quote}
That teenage boy is very clever.
\end{quote}

A masculine noun.

\begin{cousinweb}
\textbf{\gu{કિશોર}} (\emph{kiśor}) is a tatsama, Sanskrit \emph{kiśora}, which Monier-Williams gives first as \textquotedblleft{}a colt\textquotedblright{} and then as a youth.
\end{cousinweb}

\begin{grammarlens}[title={What you've built}]
One more word for describing what you lived and felt.
\end{grammarlens}

\subsection*{Your turn}
\begin{itemize}
\raggedright
  \item \emph{Say it:} \emph{kiśor}
  \item \emph{Say it:} \emph{kiśor}, once more
  \item \emph{Say it:} \emph{te kiśor ghaṇo hośiyār che.} — That teenage boy is very clever.
  \item \emph{Recall:} say the Gujarati for childhood, then the Gujarati for to grow up, then say \emph{kiśor} again
\end{itemize}

\begin{tcolorbox}[breakable,colback=teal!4,colframe=teal!35!black,title={Before you move on}]
What does \gu{કિશોર} mean? (\textquotedblleft{}A teenage boy\textquotedblright{}.) Say it once more.
\end{tcolorbox}
\section[\texorpdfstring{\gu{વૃદ્ધ} (vṛddh)}{વૃદ્ધ (vṛddh)}]{\gu{વૃદ્ધ} (vṛddh) — elderly; an old person}

\label{lesson:GU-C494-vrddh}



\begin{quote}
Before the new one: say the Gujarati for to grow up, then the Gujarati for a teenage boy.
\end{quote}

\subsection*{You'll want to know: \gu{વૃદ્ધ}}
\textbf{\gu{વૃદ્ધ}} — \emph{vṛddh} — \textquotedblleft{}elderly; an old person\textquotedblright{}.

\begin{quote}
\gu{બસમાં} \gu{વૃદ્ધ} \gu{માણસને} \gu{બેસવાની} \gu{જગ્યા} \gu{આપો}.
\end{quote}

\begin{quote}
\emph{basmā̃ vṛddh māṇasne besvānī jagyā āpo.}
\end{quote}

\begin{quote}
On the bus, give your seat to an elderly person.
\end{quote}

A masculine or feminine noun, by the person meant.

\begin{cousinweb}
\textbf{\gu{વૃદ્ધ}} (\emph{vṛddh}) is a tatsama, Sanskrit \emph{vṛddha}, \textquotedblleft{}grown, old\textquotedblright{}, as Monier-Williams records it.
\end{cousinweb}

\begin{grammarlens}[title={What you've built}]
One more word for describing what you lived and felt.
\end{grammarlens}

\subsection*{Your turn}
\begin{itemize}
\raggedright
  \item \emph{Say it:} \emph{vṛddh}
  \item \emph{Say it:} \emph{vṛddh}, once more
  \item \emph{Say it:} \emph{basmā̃ vṛddh māṇasne besvānī jagyā āpo.} — On the bus, give your seat to an elderly person.
  \item \emph{Recall:} say the Gujarati for to grow up, then the Gujarati for a teenage boy, then say \emph{vṛddh} again
\end{itemize}

\begin{tcolorbox}[breakable,colback=teal!4,colframe=teal!35!black,title={Before you move on}]
What does \gu{વૃદ્ધ} mean? (\textquotedblleft{}Elderly; an old person\textquotedblright{}.) Say it once more.
\end{tcolorbox}
\section[\texorpdfstring{\gu{વૃદ્ધાશ્રમ} (vṛddhāśram)}{વૃદ્ધાશ્રમ (vṛddhāśram)}]{\gu{વૃદ્ધાશ્રમ} (vṛddhāśram) — an old people's home}

\label{lesson:GU-C494-vrddhasram}



\begin{quote}
Before the new one: say the Gujarati for a teenage boy, then the Gujarati for elderly.
\end{quote}

\subsection*{You'll want to know: \gu{વૃદ્ધાશ્રમ}}
\textbf{\gu{વૃદ્ધાશ્રમ}} — \emph{vṛddhāśram} — \textquotedblleft{}an old people's home\textquotedblright{}.

\begin{quote}
\gu{મીરા} \gu{રવિવારે} \gu{વૃદ્ધાશ્રમમાં} \gu{ગીતો} \gu{ગાવા} \gu{જાય} \gu{છે}.
\end{quote}

\begin{quote}
\emph{Mīrā ravivāre vṛddhāśrammā̃ gīto gāvā jāy che.}
\end{quote}

\begin{quote}
Every Sunday Mira goes to sing at the old people's home.
\end{quote}

A masculine noun.

\begin{cousinweb}
\textbf{\gu{વૃદ્ધાશ્રમ}} (\emph{vṛddhāśram}) is built from \textbf{\gu{વૃદ્ધ}} (\emph{vṛddh}), \textquotedblleft{}old\textquotedblright{}, from the entry before, and \emph{āśram}, Sanskrit \emph{āśrama}, \textquotedblleft{}a hermitage, a place of rest\textquotedblright{} (Monier-Williams).
\end{cousinweb}

\begin{grammarlens}[title={What you've built}]
One more word for describing what you lived and felt.
\end{grammarlens}

\subsection*{Your turn}
\begin{itemize}
\raggedright
  \item \emph{Say it:} \emph{vṛddhāśram}
  \item \emph{Say it:} \emph{vṛddhāśram}, once more
  \item \emph{Say it:} \emph{Mīrā ravivāre vṛddhāśrammā̃ gīto gāvā jāy che.} — Every Sunday Mira goes to sing at the old people's home.
  \item \emph{Recall:} say the Gujarati for a teenage boy, then the Gujarati for elderly, then say \emph{vṛddhāśram} again
\end{itemize}

\begin{tcolorbox}[breakable,colback=teal!4,colframe=teal!35!black,title={Before you move on}]
What does \gu{વૃદ્ધાશ્રમ} mean? (\textquotedblleft{}An old people's home\textquotedblright{}.) Say it once more.
\end{tcolorbox}
